\documentclass{article}
\usepackage{amsmath}
\usepackage{xcolor}
\begin{document}

\textbf{INTR. A LAS ECUACIONES DIFERENCIALES}

Concepto. Una ecuacion diferencial es una expresion matematica 

que contempla dentro de si la presencia de derivadas y funciones.

Comunmente las funciones representan cantidades fisicas en general,

las derivadas representan la variacion o cambio de la variable 

independiente a lo largo del tiempo, y la E.D. como tal  define la relacion entre ambas.




\begin{gather*}
\text{Ecuacion } \rightarrow   x^{2}+5x-3 = 0   \rightarrow \text{ determinar el valor de x}_{1,2} = ?? \\
Por decir, si asumimos al azar que el valor de 'x'  fuese x = 2 \\
\vspace{0.5em} \\
(2)^{2}+5(2)-3 = 0 \\
\vspace{0.5em} \\
x+y=2 \\
3y-x = 5    \rightarrow  Determinar el valor x,y = ?? \\
\vspace{0.5em} \\
\frac{d^{2}y}{dx^{2}}+ 5y =e^{5x}   \rightarrow \text{ variable dependiente 'y' depende de 'x'} \\
y'' + 5y = e^{5x}   \rightarrow \text{ El objetivo es determinar el valor la } \\
\text{funcion dependiente que satisfaga la E.D. } \rightarrow  y(x) = ? \\
\vspace{0.5em} \\
\text{supongamos y}(x) = 3e^{5x}   \\
\vspace{0.5em} \\
( 3e^{5x} )''+5( 3e^{5x} ) = e^{5x} 
\end{gather*}


\textbf{Clasificacion de las E.D. de acuerdo a la derivada presente}

\colorbox[RGB]{248,231,28}{\textbf{E.D. Ordinarias}} $\rightarrow$  son aquellas que no tienen presente dentro si, derivadas parciales

\textbf{E.D. No ordinarias} $\rightarrow$  son aquellas que tienen presente derivadas paraciales



\textbf{Orden, Grado y Linealidad de una E.D. Ordinaria}

\textbf{Orden} $\rightarrow$  se representa por la derivada mayor presente en la E.D.


\begin{gather*}
\frac{d^{2}y}{dx^{2}} +\frac{dy}{dx}+y=0
\end{gather*}
             
\begin{gather*}
\frac{d^{3}y}{dx^{3}} +\frac{d^{4}y}{dx^{4}}-5y=e^{5x}
\end{gather*}
        
\begin{gather*}
y'+5xy=3x
\end{gather*}


Es. E.D 2do Orden			Es. E.D. de 4to Orden	 Es. de Primer Orden



\textbf{Grado} $\rightarrow$  el exponente mayor presente en las diferenciales, considerando 

el orden como factor de arranque.


\begin{gather*}
(\frac{d^{3}y}{dx^{3}})^{4} +(\frac{d^{4}y}{dx^{4}})^{2}+\text{\colorbox[RGB]{248,231,28}{$($}}\text{\colorbox[RGB]{248,231,28}{$\frac{d^{4}y}{dx^{4}}$}}\text{\colorbox[RGB]{248,231,28}{$)$}}\text{\colorbox[RGB]{248,231,28}{$^{3}$}}-5y=e^{5x}
\end{gather*}
  $\rightarrow$  el grado es 3





\textbf{Linealidad} $\rightarrow$  Asociado al formato o regla que deben cumplir los 

elementos presentes en una E.D.

	Ed de 2do Orden $\rightarrow$    
\begin{gather*}
Ax^{2}+Bx+C = 0
\end{gather*}
  $\rightarrow$     
\begin{gather*}
Ax^{3}+Bx^{2}+Cx+D = 0^{}
\end{gather*}


	
\begin{gather*}
5x^{2}+10 = 0
\end{gather*}
       
\begin{gather*}
5x^{2}=0
\end{gather*}
     
\begin{gather*}
5x^{2}-1x+1=0
\end{gather*}
   
\begin{gather*}
5\sqrt{x^{2}}-1x+1=0
\end{gather*}





\begin{gather*}
a_{n}(x)\frac{d^{n}y}{dx^{n}}+a_{n-1}(x)\frac{d^{n-1}y}{dx^{n-1}}+......a_{0}(x)y = g(x)
\end{gather*}



\begin{gather*}
5x\frac{d^{2}y}{dx^{2}}+10sen(x)\frac{dy}{dx} -1 y = \sqrt{y}
\end{gather*}









\end{document}
